%%=============================================================================
%% Inleiding
%%=============================================================================

\chapter{\IfLanguageName{dutch}{Inleiding}{Introduction}}%
\label{ch:inleiding}

Tijdens de Gentse Feesten verplaatsen grote groepen bezoekers zich voortdurend tussen pleinen, straten en podia. De drukte is daardoor niet gelijkmatig over de binnenstad verdeeld: een locatie kan in korte tijd vollopen, terwijl het elders relatief rustig blijft. Voor hulpdiensten is het daarom niet alleen relevant hoeveel bezoekers op het evenement aanwezig zijn, maar vooral waar de drukte toeneemt, hoe snel dat gebeurt en welke doorgangen daarbij onder druk komen te staan. Recente veiligheidsmaatregelen, zoals het spreiden van de uitstroom na een bijzonder drukke feestendag, illustreren de nood aan een actueel beeld van bezoekersstromen \autocite{Aneca2024}.

Camera's kunnen lokale veranderingen continu zichtbaar maken, maar het gelijktijdig beoordelen van meerdere videobeelden blijft arbeidsintensief. Computervisie biedt een manier om die beelden automatisch te analyseren. In zeer drukke scènes is het vaak moeilijk om elke persoon afzonderlijk te detecteren, bijvoorbeeld door overlap, perspectiefverschillen en wisselende belichting. Daarom onderzoekt deze bachelorproef \emph{density estimation}: een techniek die een dichtheidskaart voorspelt waarvan de som een schatting geeft van het aantal personen in beeld \autocite{Khan2020}. Zo'n kaart kan drukke zones zichtbaar maken en hulpdiensten van aanvullende, objectieve informatie voorzien.

\section{\IfLanguageName{dutch}{Probleemstelling}{Problem Statement}}%
\label{sec:probleemstelling}

De Gentse Feesten vinden sinds 1843 plaats en groeiden uit tot een tiendaags cultureel volksfeest in de historische binnenstad. In 2023 werden meer dan 1,6 miljoen bezoekers geteld, verspreid over dag en nacht \autocite{StadGentGeschiedenis}. Omdat het evenement niet op één afgesloten terrein plaatsvindt, verplaatsen bezoekers zich vrij tussen locaties. Een algemeen bezoekersaantal zegt daardoor weinig over lokale drukte of mogelijke knelpunten.

Een hoge personendichtheid kan de bewegingsvrijheid beperken en het risico op gevaarlijke situaties verhogen. Organisatoren en hulpdiensten baseren zich daarvoor onder meer op terreinobservaties, handmatige tellingen en ervaring uit voorgaande edities. Die werkwijze blijft onmisbaar, maar kan niet gelijktijdig elke locatie opvolgen en leidt onvermijdelijk tot een vertraging tussen waarnemen, afstemmen en handelen.

Stad Gent experimenteerde reeds met meerdere vormen van druktemeting. Het project VLOED wijst op de nood aan fijnmazige, actuele informatie over de ruimtelijke spreiding van drukte \autocite{StadGentVloed}. Er is echter een bijkomende vertaalslag nodig om camerabeelden om te zetten in informatie die operationeel bruikbaar is. Dit onderzoek verkent in welke mate een density-estimationmodel daarvoor een technische basis kan vormen.

\section{\IfLanguageName{dutch}{Onderzoeksvraag}{Research question}}%
\label{sec:onderzoeksvraag}

Hoe kan computervisie, meer bepaald density estimation, bruikbare en tijdige indicaties van lokale drukte leveren tijdens openluchtevenementen zoals de Gentse Feesten?

\textit{Deelvragen rond het probleemdomein:}
\begin{itemize}
    \item Welke factoren maken lokale drukte risicovol, en welke informatie hebben hulpdiensten nodig om proactief te kunnen ingrijpen?
    \item Welke bestaande technieken voor druktemeting zijn beschikbaar en welke beperkingen hebben ze voor een verspreid stadsfestival?
\end{itemize}

\textit{Deelvragen rond het oplossingsdomein:}
\begin{itemize}
    \item Waarom is density estimation een geschikte aanvulling op persoonsdetectie in dicht bezette scènes?
    \item Welke dataset, evaluatiemetrieken en experimentele opzet zijn nodig om een eerste model reproduceerbaar te beoordelen?
    \item Welke nauwkeurigheid en inferentiesnelheid bereikt een CSRNet-baseline op ShanghaiTech Part B?
    \item Welke bijkomende kalibratie is nodig om een dichtheidskaart in beeldcoördinaten te vertalen naar een fysieke dichtheid in personen per vierkante meter?
\end{itemize}

\section{\IfLanguageName{dutch}{Onderzoeksdoelstelling}{Research objective}}%
\label{sec:onderzoeksdoelstelling}

Het doel van dit onderzoek is een reproduceerbare Proof of Concept (PoC) te realiseren die met CSRNet een density map en een geschat aantal personen uit één camerabeeld afleidt. De PoC wordt geëvalueerd op ShanghaiTech Part B aan de hand van de Mean Absolute Error (MAE), de Root Mean Squared Error (RMSE) en de gemeten inferentietijd. De resultaten tonen welke prestaties een eerste baseline haalt en welke stappen nog nodig zijn voor inzet in de context van de Gentse Feesten.

De PoC is beslissingsondersteunend: zij kan een operator attenderen op relatieve veranderingen of drukke zones, maar vervangt geen menselijke beoordeling. De dichtheidskaart drukt in de huidige opzet een verdeling in beeldpixels uit. Ze kan dus niet zonder camerakalibratie worden geïnterpreteerd als een fysieke dichtheid in personen per vierkante meter, en vormt op zichzelf geen grond voor een veiligheidsalarm. Het onderzoek richt zich evenmin op het herkennen of volgen van individuele bezoekers.

\section{\IfLanguageName{dutch}{Opzet van deze bachelorproef}{Structure of this bachelor thesis}}%
\label{sec:opzet-bachelorproef}

Hoofdstuk~\ref{ch:stand-van-zaken} situeert de Gentse Feesten als casus, bespreekt risico's en bestaande technieken voor druktemeting en motiveert de keuze voor density estimation. Hoofdstuk~\ref{ch:methodologie} beschrijft de onderzoeksmethode, de datasetsplitsing en het evaluatieprotocol.

Hoofdstuk~\ref{ch:poc} licht de implementatie van CSRNet toe en rapporteert de resultaten van de uitgevoerde CPU- en GPU-experimenten. Ten slotte beantwoordt Hoofdstuk~\ref{ch:conclusie} de onderzoeksvraag, plaatst het de resultaten in hun context en formuleert het aanbevelingen voor verder onderzoek en eventuele operationele inzet.
